\subsection{Can the boundary be outrun by refining the unit?}\label{sec:domains}

A practitioner facing~\eqref{eq:boundary} has one obvious move: refine the unit.
Splitting countries into regions or demographic domains multiplies the
calibration count. Refinement is not free, though. Smaller domains carry smaller
samples, the design variance rises, the design share rises with it, and so does
the requirement. Which side wins decides whether the boundary can be engineered
around.

We hold the estimand fixed as a domain's departure from its own national
distribution and vary only how domains are cut, over three rounds and three
minimum sizes. Table~\ref{tab:refine} gives the result, and the answer is that
the number of domains is not what matters.

\begin{table}[htbp]\centering
\caption{Ways to cut the same survey into domains, at a minimum domain size of
60, averaged over rounds. ``Required'' is~\eqref{eq:boundary} evaluated at the
design share each cut produces.}\label{tab:refine}
\begin{tabular}{lrrrr}\toprule
domains cut by & $K$ & design share & required & shortfall\\\midrule
sex & \DomSexK & \DomSexRho & \DomSexReq & $-1946$\\
age band & \DomAgeK & \DomAgeRho & \DomAgeReq & $-70$\\
region & \DomRegK & \DomRegRho & \DomRegReq & $-20$\\
age band and sex & \DomAgeSexK & \DomAgeSexRho & \DomAgeSexReq & $-98$\\
region and sex & \DomRegSexK & \DomRegSexRho & \DomRegSexReq & $-101$\\
\bottomrule\end{tabular}
\end{table}

\paragraph{The dimension matters more than the count.} Splitting by sex yields
\DomSexK{} populations against a requirement of \DomSexReq. Splitting by region
yields \DomRegK{} against \DomRegReq, within reach. The difference is not the
number of domains, which differs only fourfold, but the design share: sex
produces \DomSexRho{} and region \DomRegRho. Halving each sample by sex adds
sampling variance without adding between-domain dispersion, because men and
women have nearly the same trust distribution. Cutting by region adds both,
because regions genuinely differ.

Adding a second, homogeneous dimension on top behaves the same way in miniature.
Region and sex together give more populations than region alone, but the
requirement rises faster still, and the cut is worse than the one it refines.

The design instruction follows, and it is the useful thing the boundary yields:
\textbf{refine along a dimension on which the domains actually differ.} A cut
that only subdivides a homogeneous population raises the requirement far faster
than the supply, and no amount of it reaches the correction. A cut that tracks
real heterogeneity can.

\paragraph{The boundary holds across domain definitions.} Over \DomConfigs{}
configurations spanning \DomTypes{} domain definitions, three rounds and three
size thresholds, \eqref{eq:boundary} predicts the reliability gate correctly in
\DomCorrect. The gate opens in \DomOpen, \DomRoundEleven{} of them in round~11,
the round with the lowest design share throughout. Those are the configurations
in which the correction may be used on this survey, and
Section~\ref{sec:widths} measures what it is worth there.
